% Zumdahl Chapter 18 -- Electrochemistry. Elaborated notes, 5 blocks.
% Covers CED 9.8-9.11, the back half of Unit 9 -- and the last of the CED.
\documentclass{apchem-notes}

\apcunitword{Chapter}
\apcunit{18}
\apcunittitle{Electrochemistry}
\apcdoctitle{Guided Notes}
\apcsubtitle{Zumdahl \S18.1--18.4, 18.7--18.8 \textbullet\ PDF pp.~881--930 \textbullet\ 5 blocks}

\begin{document}
\apctitleblock

\begin{apcnote}[How this chapter fits the AP course]
\textbf{This is the last chapter of the course, and it completes the CED.}
\S18.1 $\to$ CED 9.8, \S18.2--18.3 $\to$ 9.9, \S18.4 $\to$ 9.10,
\S18.7--18.8 $\to$ 9.11. With Chapter 17 it finishes Unit 9, and with it the
entire framework.

\textbf{Skip:} \S18.5 (batteries) and \S18.6 (corrosion) are applications,
not framework topics --- interesting, and referenced in passing, but not
examined.

Electrochemistry is where three earlier units meet. Redox bookkeeping comes
from \textbf{Unit 4}, the $Q$-versus-$K$ argument from \textbf{Unit 7}, and
$\Delta G^\circ$ from \textbf{Chapter 17}. The single new idea is that a
favored redox reaction can be made to push electrons through a wire, and
that the push is measurable as a voltage.
\end{apcnote}

\begin{apctrap}
\textbf{Two CED boundaries that change what you should study.}

\textbf{1. Do not label electrodes ``positive'' or ``negative.''} Topic 9.8
carries an explicit exclusion: \emph{labeling an electrode as positive or
negative will not be assessed on the AP Exam.} The sign convention differs
between galvanic and electrolytic cells and is a reliable source of
confusion, so the exam simply avoids it. Learn \textbf{anode} and
\textbf{cathode} instead --- those never change meaning.

\textbf{2. The Nernst equation is not required.} Topic 9.10 asks you to
\emph{reason} about how a cell potential changes when conditions move away
from standard --- using $Q$ against $K$ --- and it lists no equation.
Zumdahl \S18.4 works through the full Nernst calculation; that is beyond the
framework. Block 4 teaches the reasoning the CED actually asks for, and
flags the algebra as enrichment.
\end{apctrap}

% =====================================================================
\blocklesson{Galvanic Cells: How They Work}{Zumdahl \S18.1}

\begin{objectives}
  \item Identify each component of an electrochemical cell and its role.
  \item Determine the direction of electron and ion flow.
\end{objectives}

\reading{Zumdahl \S18.1}{882--886}

\begin{warmup}[4]
  \item Oxidation is the \answerblank[3cm]{loss} of electrons.
  \item Reduction is the \answerblank[3cm]{gain} of electrons.
  \item Oxidation number of \ce{Zn} in \ce{Zn(s)}:
        \answerblank[1.6cm]{0}
  \item A reaction with $\Delta G^\circ < 0$ is:
        \answerblank[5.4cm]{thermodynamically favored}
\end{warmup}

\segment{INSTRUCTION A}{25}{Separating the half-reactions}

\notesectionz{Why a wire is involved at all}{18.1}
\practice{2}

Drop zinc metal into copper(II) sulfate and the reaction
\[ \ce{Zn(s) + Cu^2+(aq) -> Zn^2+(aq) + Cu(s)} \]
happens immediately, releasing energy as \answerblank[1.6cm]{heat}. The
electrons pass directly from zinc atoms to copper ions at the point of
contact, and nothing useful is captured.

A \term{galvanic cell} (also called a voltaic cell) does the same chemistry
with the two halves \textbf{physically separated}, so the electrons are
forced to travel through an external wire to get from one to the other.
That flow is an electric current, and it can do work.

\notesectionz{The parts, and what each one does}{18.1}

\begin{center}
\begin{tabular}{@{}>{\raggedright\arraybackslash}p{3.2cm}%
                 >{\raggedright\arraybackslash}p{9.6cm}@{}}
\hline
\textbf{Component} & \textbf{Role} \\
\hline
Anode & where \textbf{oxidation} happens; electrons are released here \\
Cathode & where \textbf{reduction} happens; electrons are consumed here \\
Salt bridge & lets ions migrate to keep each half-cell electrically
  neutral; without it the cell stops almost immediately \\
External wire & carries electrons from anode to cathode \\
Voltmeter & measures the potential difference driving that flow \\
\hline
\end{tabular}
\end{center}

\begin{center}
\fbox{\parbox{0.8\linewidth}{\centering
\textbf{An Ox} --- \textbf{An}ode, \textbf{Ox}idation \\
\textbf{Red Cat} --- \textbf{Red}uction, \textbf{Cat}hode \\[0.3em]
True in \emph{every} cell, galvanic and electrolytic alike.}}
\end{center}

Electrons always flow from \answerblank[2.4cm]{anode} to
\answerblank[2.4cm]{cathode} through the wire --- alphabetical order, which
is worth noticing because it never changes.

\begin{apcnote}[What the salt bridge is actually for]
As the cell runs, the anode compartment accumulates positive ions
(\ce{Zn^2+} entering solution) and the cathode compartment loses them
(\ce{Cu^2+} plating out). Charge would build up on both sides within
moments and the electron flow would stop.

The salt bridge prevents that by letting spectator ions migrate: anions move
\emph{toward the anode} and cations \emph{toward the cathode}, cancelling
the charge imbalance. It completes the circuit without letting the two
solutions mix.
\end{apcnote}

\segment{GUIDED PRACTICE}{15}{Reading a cell}

For $\ce{Zn(s) + Cu^2+(aq) -> Zn^2+(aq) + Cu(s)}$ run as a galvanic cell:

\begin{enumerate}[leftmargin=*,itemsep=0.35em]
  \item Oxidation half-reaction:
        \answerblank[6cm]{\ce{Zn -> Zn^2+ + 2e^-}}
  \item Reduction half-reaction:
        \answerblank[6cm]{\ce{Cu^2+ + 2e^- -> Cu}}
  \item The anode is made of: \answerblank[3cm]{zinc}
  \item The cathode is made of: \answerblank[3cm]{copper}
  \item Electrons flow from \answerblank[5.4cm]{the zinc to the copper}
  \item The zinc electrode's mass \answerblank[2.4cm]{decreases}; the
        copper electrode's mass \answerblank[2.4cm]{increases}
  \item Anions in the salt bridge move toward:
        \answerblank[3.6cm]{the anode (zinc)}
\end{enumerate}

\segment{INSTRUCTION B}{20}{Galvanic against electrolytic}

\notesectionz{The same hardware, run backwards}{18.7}
\practice{2}

\begin{center}
\begin{tabular}{@{}lll@{}}
\hline
 & \textbf{Galvanic (voltaic)} & \textbf{Electrolytic} \\
\hline
Reaction & thermodynamically \textbf{favored} & thermodynamically
  \textbf{unfavored} \\
$E^\circ_{\text{cell}}$ & positive & negative \\
$\Delta G^\circ$ & negative & positive \\
Energy & \emph{produces} electrical energy & \emph{consumes} it from an
  external source \\
Anode & oxidation & oxidation \\
Cathode & reduction & reduction \\
\hline
\end{tabular}
\end{center}

The bottom two rows are the point of the table: \textbf{anode and cathode
mean the same thing in both}. What reverses is which direction the reaction
is being pushed, and whether energy comes out or must be put in.

\begin{apctrap}
This is exactly Chapter 17's ``driving an unfavorable process with an
external energy source,'' in hardware. An electrolytic cell is CED topic
9.7 made physical --- and it is why electrolysis was one of the CED's two
named examples of external energy.
\end{apctrap}

\segment{APPLICATION}{20}{Describing cells}

\begin{enumerate}[leftmargin=*,itemsep=0.4em]
  \item A student builds a cell and finds the silver electrode gains mass
        while the copper electrode dissolves. Identify the anode, the
        cathode, and the direction of electron flow.
        \answerlines[3]{Copper is the anode --- it is being oxidized and
        losing mass. Silver is the cathode --- \ce{Ag+} is being reduced
        and plating onto it. Electrons flow through the wire from the
        copper to the silver.}
  \item Why does a galvanic cell stop working if the salt bridge is
        removed? \answerlines[3]{Charge builds up: the anode compartment
        becomes positively charged as cations enter solution, and the
        cathode compartment becomes negatively charged as cations are
        removed. That imbalance opposes further electron flow, and the
        current halts almost at once.}
  \item A cell has $E^\circ_{\text{cell}} = -0.46$~V. Is it galvanic or
        electrolytic, and what must be supplied?
        \answerlines[2]{Electrolytic. A negative cell potential means
        $\Delta G^\circ > 0$, so the reaction is unfavored and an external
        power supply must apply a potential greater than 0.46 V to drive
        it.}
\end{enumerate}

\begin{exitticket}
State the one rule about anodes and cathodes that is true in every
electrochemical cell, and one thing the exam will \emph{not} ask you to
label.
\keyonly{\begin{apcanswer}Oxidation always occurs at the anode and reduction
always at the cathode --- in galvanic and electrolytic cells alike. The exam
will not ask you to label an electrode as positive or negative; that
convention differs between the two cell types and is explicitly
excluded.\end{apcanswer}}
\end{exitticket}

% =====================================================================
\blocklesson{Standard Reduction Potentials}{Zumdahl \S18.2}

\begin{objectives}
  \item Calculate $E^\circ_{\text{cell}}$ from standard reduction
        potentials.
  \item Predict which species is oxidized and which is reduced.
\end{objectives}

\reading{Zumdahl \S18.2}{886--892}

\begin{warmup}[3]
  \item Oxidation occurs at the: \answerblank[2.4cm]{anode}
  \item Electrons flow from anode to:
        \answerblank[2.4cm]{cathode}
  \item A favored reaction has $E^\circ_{\text{cell}}$:
        \answerblank[2.4cm]{positive}
\end{warmup}

\segment{INSTRUCTION A}{25}{A scale built on an arbitrary zero}

\notesectionz{The standard hydrogen electrode}{18.2}
\practice{5}

Only a \emph{difference} in potential can be measured --- there is no way to
measure one half-cell alone. So one half-reaction is assigned a value by
convention:
\[ \ce{2H+(aq) + 2e^- -> H2(g)} \qquad E^\circ \equiv
   \mathbf{0.00}~\text{V} \]
Every other half-reaction is then measured against this
\term{standard hydrogen electrode}. Zumdahl's zinc cell reads 0.76 V against
it, so \ce{Zn^2+ + 2e^- -> Zn} is assigned $E^\circ = -0.76$~V.

\textbf{Standard conditions} means all solutes at \SI{1}{\Molar}, all gases
at \SI{1}{\atm}, and \SI{25}{\celsius}.

\begin{apctrap}
\textbf{Every table is written as reductions.} That is a universal
convention, and it means the table does not tell you which way a reaction
runs --- you decide that.

Reversing a half-reaction \textbf{flips the sign} of $E^\circ$. But
multiplying a half-reaction by a coefficient does \textbf{not} change
$E^\circ$ at all. Potential is an intensive property --- volts per charge,
not volts per mole. Scaling a half-reaction to balance electrons is
required, and scaling its voltage is wrong.
\end{apctrap}

\notesectionz{Computing the cell potential}{18.2}

\[ \boxed{\;E^\circ_{\text{cell}} = E^\circ_{\text{cathode}}
   - E^\circ_{\text{anode}}\;} \]
where both values are taken \emph{straight from the table as reductions} ---
no sign flipping needed if you use this form.

For a \textbf{galvanic} cell you want $E^\circ_{\text{cell}} > 0$, so the
species with the \answerblank[3cm]{more positive} $E^\circ$ is reduced (it
is the cathode) and the other is oxidized.

\begin{workedexample}[Worked example 1: the Daniell cell]
Combine \ce{Zn^2+/Zn} ($-0.76$~V) and \ce{Cu^2+/Cu} ($+0.34$~V).

Copper has the more positive reduction potential, so \textbf{copper is
reduced} (cathode) and \textbf{zinc is oxidized} (anode):
\begin{align*}
  \text{cathode:} &\quad \ce{Cu^2+ + 2e^- -> Cu} \\
  \text{anode:}   &\quad \ce{Zn -> Zn^2+ + 2e^-} \\
  \text{overall:} &\quad \ce{Zn + Cu^2+ -> Zn^2+ + Cu}
\end{align*}
\[ E^\circ_{\text{cell}} = 0.34 - (-0.76) = \mathbf{+1.10}~\text{V} \]
Positive, so the cell is galvanic --- and 1.10 V is exactly what a
voltmeter reads on this cell.
\end{workedexample}

\begin{workedexample}[Worked example 2: when the electrons do not balance]
Combine \ce{Ag+/Ag} ($+0.80$~V) and \ce{Cu^2+/Cu} ($+0.34$~V).

Silver is more positive, so silver is reduced and copper oxidized. Balancing
electrons requires \emph{two} silver half-reactions:
\begin{align*}
  \text{cathode:} &\quad \ce{2Ag+ + 2e^- -> 2Ag} \\
  \text{anode:}   &\quad \ce{Cu -> Cu^2+ + 2e^-} \\
  \text{overall:} &\quad \ce{Cu + 2Ag+ -> Cu^2+ + 2Ag}
\end{align*}
\[ E^\circ_{\text{cell}} = 0.80 - 0.34 = \mathbf{+0.46}~\text{V} \]

\textbf{Note what did not happen:} the silver half-reaction was doubled, but
its $E^\circ$ was \emph{not}. Writing $2(0.80) - 0.34 = 1.26$~V is the
classic error in this chapter.
\end{workedexample}

\segment{GUIDED PRACTICE}{15}{Building cells from the table}

$E^\circ$ (V): \ce{Ag+/Ag} $+0.80$ \textbullet\ \ce{Cu^2+/Cu} $+0.34$
\textbullet\ \ce{Pb^2+/Pb} $-0.13$ \textbullet\ \ce{Ni^2+/Ni} $-0.23$
\textbullet\ \ce{Fe^2+/Fe} $-0.44$ \textbullet\ \ce{Zn^2+/Zn} $-0.76$
\textbullet\ \ce{Mg^2+/Mg} $-2.37$

\begin{enumerate}[leftmargin=*,itemsep=0.35em]
  \item \ce{Zn} and \ce{Ag+}: $E^\circ_{\text{cell}} = $
        \answerblank[2.6cm]{$+1.56$ V}
  \item \ce{Mg} and \ce{Cu^2+}: $E^\circ_{\text{cell}} = $
        \answerblank[2.6cm]{$+2.71$ V}
  \item \ce{Zn} and \ce{Ni^2+}: $E^\circ_{\text{cell}} = $
        \answerblank[2.6cm]{$+0.53$ V}
  \item \ce{Fe} and \ce{Cu^2+}: $E^\circ_{\text{cell}} = $
        \answerblank[2.6cm]{$+0.78$ V}
  \item \ce{Pb} and \ce{Ag+}: $E^\circ_{\text{cell}} = $
        \answerblank[2.6cm]{$+0.93$ V}
\end{enumerate}

\segment{APPLICATION}{20}{Reading the table as a ranking}

\begin{enumerate}[leftmargin=*,itemsep=0.4em]
  \item Which is the stronger oxidizing agent, \ce{Ag+} or \ce{Zn^2+}?
        Justify. \answerlines[2]{\ce{Ag+}. A more positive reduction
        potential means the species is more readily reduced, and a species
        that is readily reduced is by definition a strong oxidizing agent.}
  \item Which is the stronger reducing agent, \ce{Mg} or \ce{Cu}?
        \answerlines[2]{\ce{Mg}. Its reduction potential is far more
        negative, so the reverse process --- oxidation of the metal --- is
        strongly favored, which is what a good reducing agent does.}
  \item Will \ce{Cu(s)} react with \SI{1}{\Molar} \ce{HCl}? Support your
        answer with a cell potential. \workspace[2.4cm]
        \keyonly{\begin{apcanswer}The proposed reaction is
        $\ce{Cu + 2H+ -> Cu^2+ + H2}$. Cathode \ce{H+/H2} $= 0.00$, anode
        \ce{Cu^2+/Cu} $= +0.34$:
        $E^\circ = 0.00 - 0.34 = \mathbf{-0.34}$~V. Negative, so
        \textbf{no} --- copper does not dissolve in hydrochloric
        acid.\end{apcanswer}}
\end{enumerate}

\begin{exitticket}
A student doubles a half-reaction to balance electrons and doubles its
$E^\circ$ as well. Explain why the second step is wrong.
\keyonly{\begin{apcanswer}Potential is an intensive property --- energy per
unit charge --- so it does not scale with the amount of substance. Doubling
the half-reaction doubles both the energy and the charge, leaving the volts
unchanged. Only the number of electrons $n$ changes, and that enters later
through $\Delta G^\circ = -nFE^\circ$.\end{apcanswer}}
\end{exitticket}

% =====================================================================
\blocklesson{Cell Potential, Free Energy, and $K$}{Zumdahl \S18.3}

\begin{objectives}
  \item Relate $E^\circ_{\text{cell}}$ to $\Delta G^\circ$ and to $K$.
  \item Explain why $n$ matters for free energy but not for voltage.
\end{objectives}

\reading{Zumdahl \S18.3}{892--895}

\begin{warmup}[3]
  \item $E^\circ_{\text{cell}} = $
        \answerblank[5cm]{$E^\circ_{\text{cathode}} -
        E^\circ_{\text{anode}}$}
  \item Reversing a half-reaction does what to $E^\circ$?
        \answerblank[3cm]{flips its sign}
  \item Doubling a half-reaction does what to $E^\circ$?
        \answerblank[2.6cm]{nothing}
\end{warmup}

\segment{INSTRUCTION A}{25}{Connecting volts to kilojoules}

\notesectionz{$\Delta G^\circ = -nFE^\circ$}{18.3}
\practice{5}

Electrical work is charge moved through a potential difference. For $n$
moles of electrons carrying \term{Faraday's constant}
$F = \SI{96485}{\coulomb\per\mole}$ of charge each,
\[ \boxed{\;\Delta G^\circ = -nFE^\circ_{\text{cell}}\;} \]

Read off the consequences:

\begin{itemize}[leftmargin=*,itemsep=0.3em]
  \item $E^\circ > 0 \Rightarrow \Delta G^\circ$
        \answerblank[2cm]{negative} $\Rightarrow$ favored, galvanic.
  \item $E^\circ < 0 \Rightarrow \Delta G^\circ$
        \answerblank[2cm]{positive} $\Rightarrow$ unfavored, electrolytic.
  \item The negative sign is what makes ``positive voltage'' and
        ``favored'' line up.
\end{itemize}

\begin{apctrap}
\textbf{Here is where $n$ finally matters.} Voltage does not depend on how
many electrons are transferred, but free energy does --- because
$\Delta G^\circ$ is an amount of energy, and more electrons moving through
the same potential means more energy.

So two cells can have the same $E^\circ$ and very different
$\Delta G^\circ$. Getting $n$ wrong is the most common way to lose a point
on this calculation, and $n$ comes from the \emph{balanced overall
equation}, not from either half-reaction alone.
\end{apctrap}

\begin{workedexample}[Worked example 3: the Daniell cell in kilojoules]
$E^\circ = +1.10$~V and $n = 2$ (two electrons per zinc atom):
\[ \Delta G^\circ = -(2)(96485)(1.10) = -212{,}267~\text{J}
   = \mathbf{-212}~\text{kJ} \]
Strongly favored, as the positive voltage already told us. Note the answer
comes out in \textbf{joules} --- $F$ is in coulombs per mole and volts are
joules per coulomb --- so it must be converted to kJ.
\end{workedexample}

\segment{INSTRUCTION B}{20}{And on to the equilibrium constant}

\notesectionz{Three descriptions of one thing}{18.3}
\practice{5}

Combining $\Delta G^\circ = -nFE^\circ$ with Chapter 17's
$\Delta G^\circ = -RT\ln K$:
\[ nFE^\circ = RT\ln K \qquad\Longrightarrow\qquad
   \ln K = \frac{nFE^\circ}{RT} \]

You now have \emph{three} equivalent ways to say ``products are favored'':

\begin{center}
\begin{tabular}{@{}cccl@{}}
\hline
$E^\circ_{\text{cell}}$ & $\Delta G^\circ$ & $K$ & \\
\hline
$> 0$ & $< 0$ & $> 1$ & favored; galvanic \\
$= 0$ & $= 0$ & $= 1$ & at equilibrium; a dead battery \\
$< 0$ & $> 0$ & $< 1$ & unfavored; electrolytic \\
\hline
\end{tabular}
\end{center}

\begin{workedexample}[Worked example 4: how far does the Daniell cell go?]
With $E^\circ = 1.10$~V and $n = 2$:
\[ \ln K = \frac{(2)(96485)(1.10)}{(8.314)(298)} = 85.7
   \qquad K = e^{85.7} \approx \mathbf{2\times10^{37}} \]
Effectively complete. A modest-looking 1.10 V corresponds to an equilibrium
constant of $10^{37}$ --- because $E^\circ$ is multiplied by $nF$, a very
large number, before it enters an exponential.
\end{workedexample}

\segment{APPLICATION}{20}{Moving between the three}

\begin{enumerate}[leftmargin=*,itemsep=0.4em]
  \item For $\ce{Cu + 2Ag+ -> Cu^2+ + 2Ag}$, $E^\circ = +0.46$~V.
        Calculate $\Delta G^\circ$. \workspace[2.2cm]
        \keyonly{\begin{apcanswer}$n = 2$;
        $\Delta G^\circ = -(2)(96485)(0.46) = -88{,}766~\text{J} =
        \mathbf{-89}~\text{kJ}$\end{apcanswer}}
  \item A cell has $E^\circ = -1.10$~V with $n = 2$. Find
        $\Delta G^\circ$ and state what it means practically.
        \workspace[2.4cm]
        \keyonly{\begin{apcanswer}$\Delta G^\circ = -(2)(96485)(-1.10) =
        \mathbf{+212}~\text{kJ}$. Unfavored --- this is the reverse Daniell
        cell, and running it requires an external supply of more than
        1.10 V.\end{apcanswer}}
  \item Two cells both have $E^\circ = +0.50$~V, but one transfers 1
        electron and the other 4. Compare their voltages and their
        $\Delta G^\circ$ values. \answerlines[3]{The voltages are equal ---
        potential is intensive and does not depend on $n$. The free energy
        changes are not: the four-electron cell has four times the
        magnitude of $\Delta G^\circ$, because four times as much charge
        moves through the same potential difference.}
\end{enumerate}

\begin{exitticket}
A battery is described as ``dead.'' Give the value of $E_{\text{cell}}$, of
$\Delta G$, and the relationship between $Q$ and $K$ at that moment.
\keyonly{\begin{apcanswer}$E_{\text{cell}} = 0$, $\Delta G = 0$, and
$Q = K$. The cell has reached equilibrium: there is no longer any driving
force to push electrons through the wire. A dead battery is not empty --- it
is at equilibrium.\end{apcanswer}}
\end{exitticket}

% =====================================================================
\blocklesson{Nonstandard Conditions}{Zumdahl \S18.4}

\begin{objectives}
  \item Predict how a change in concentration changes the cell potential.
  \item Explain how a concentration cell produces a voltage.
\end{objectives}

\reading{Zumdahl \S18.4}{895--901}

\begin{warmup}[3]
  \item $\Delta G^\circ = $ \answerblank[3cm]{$-nFE^\circ$}
  \item $Q = K$ means the system is at:
        \answerblank[2.6cm]{equilibrium}
  \item Faraday's constant: \answerblank[3.4cm]{96,485 C/mol}
\end{warmup}

\begin{apcnote}[What the CED asks for here]
Topic 9.10 lists \textbf{no equation}. It asks you to \emph{reason} about
how the cell potential responds when conditions leave standard state, using
the distance from equilibrium. The Nernst equation in Zumdahl \S18.4 is
\textbf{enrichment} --- useful, but not required, and not on this chapter's
test.
\end{apcnote}

\segment{INSTRUCTION A}{25}{Voltage as distance from equilibrium}

\notesectionz{The one idea in this block}{18.4}
\practice{6}

\begin{center}
\fbox{\parbox{0.88\linewidth}{\centering
The cell potential is a \textbf{driving force toward equilibrium}. \\
The further the cell is from equilibrium, the \textbf{larger} $|E|$. \\
At equilibrium, $Q = K$ and $E = \mathbf{0}$.}}
\end{center}

Standard conditions correspond to $Q = \answerblank[1.4cm]{1}$, which is
where $E = E^\circ$. So to decide whether a change raises or lowers the
voltage, ask one question: \textbf{does the change move $Q$ further from
$K$, or closer to it?}

\begin{center}
\begin{tabular}{@{}ll@{}}
\hline
Change takes the cell \emph{further} from equilibrium & $|E|$
  \textbf{increases} \\
Change takes the cell \emph{closer} to equilibrium & $|E|$
  \textbf{decreases} \\
\hline
\end{tabular}
\end{center}

For a favored cell, $K$ is large, so $Q$ starts below it. Then:
\emph{increasing reactant} concentration lowers $Q$, moving it
\answerblank[2.4cm]{further} from $K$, so the voltage
\answerblank[2cm]{rises}. \emph{Increasing product} concentration raises
$Q$, moving it closer to $K$, so the voltage
\answerblank[2cm]{falls}.

\begin{apctrap}
\textbf{Do not invoke Le Ch\^atelier here.} The CED says so outright:
equilibrium arguments such as Le Ch\^atelier's principle \emph{do not apply
to electrochemical systems, because the systems are not at equilibrium}.

A running cell is deliberately held away from equilibrium --- that is where
its voltage comes from. Reason with $Q$ against $K$ and with distance from
equilibrium, not with ``the equilibrium shifts.''
\end{apctrap}

\begin{workedexample}[Worked example 5: which way does the voltage move?]
For the Daniell cell,
$\ce{Zn(s) + Cu^2+(aq) -> Zn^2+(aq) + Cu(s)}$, so
$Q = [\ce{Zn^2+}]/[\ce{Cu^2+}]$.

\textbf{Raise $[\ce{Cu^2+}]$.} $Q$ gets \emph{smaller}, so the cell is
further from equilibrium and $E$ \textbf{increases} above 1.10 V.

\textbf{Raise $[\ce{Zn^2+}]$.} $Q$ gets \emph{larger}, closer to $K$, so
$E$ \textbf{decreases}.

\textbf{Add more zinc metal.} No change at all --- \ce{Zn(s)} is a pure
solid and does not appear in $Q$.
\end{workedexample}

\segment{INSTRUCTION B}{20}{Concentration cells}

\notesectionz{A voltage from nothing but a difference}{18.4}
\practice{6}

A \term{concentration cell} has the \emph{same} half-reaction on both sides,
differing only in concentration. Since both electrodes are identical,
\[ E^\circ_{\text{cell}} = 0.34 - 0.34 = \mathbf{0}~\text{V} \]
and yet the cell produces a measurable voltage. The reason is that it is not
at equilibrium: equilibrium here means \textbf{equal concentrations}.

\begin{workedexample}[Worked example 6: reasoning out the direction]
\[ \ce{Cu} \mid \ce{Cu^2+}~(\SI{0.010}{\Molar}) \parallel
   \ce{Cu^2+}~(\SI{1.0}{\Molar}) \mid \ce{Cu} \]

Ask what has to happen for the two sides to become equal:

\begin{itemize}[leftmargin=*,itemsep=0.25em]
  \item The \textbf{dilute} side must become more concentrated, so
        \ce{Cu^2+} must be \emph{produced} there:
        $\ce{Cu -> Cu^2+ + 2e^-}$. That is oxidation, so the dilute side is
        the \textbf{anode}.
  \item The \textbf{concentrated} side must become more dilute, so
        \ce{Cu^2+} must be \emph{consumed}:
        $\ce{Cu^2+ + 2e^- -> Cu}$. Reduction --- the \textbf{cathode}.
\end{itemize}

Electrons therefore flow from the dilute half-cell to the concentrated one.
The voltage is small but positive, and it falls to \textbf{zero} once the
two concentrations match.
\end{workedexample}

\segment{APPLICATION}{20}{Predicting the response}

\begin{enumerate}[leftmargin=*,itemsep=0.4em]
  \item For $\ce{Zn + Cu^2+ -> Zn^2+ + Cu}$, state and justify the effect
        on $E$ of (i) diluting the \ce{Cu^2+} solution, (ii) adding
        \ce{Na2S} to the cathode compartment to precipitate \ce{CuS}.
        \answerlines[4]{(i) Lowering $[\ce{Cu^2+}]$ raises
        $Q = [\ce{Zn^2+}]/[\ce{Cu^2+}]$, moving the cell closer to
        equilibrium, so $E$ decreases. (ii) Precipitating \ce{CuS} removes
        \ce{Cu^2+} from solution, which is a further and larger decrease in
        $[\ce{Cu^2+}]$ --- $Q$ rises more, and $E$ decreases further
        still.}
  \item A concentration cell is built from \ce{Ag} electrodes in
        \SI{0.10}{\Molar} and \SI{2.0}{\Molar} \ce{AgNO3}. Identify the
        anode and predict what happens to $E$ over time.
        \answerlines[3]{The \SI{0.10}{\Molar} half-cell is the anode ---
        silver dissolves there to raise its concentration. The
        \SI{2.0}{\Molar} side is the cathode. As the cell runs the
        concentrations converge, so $E$ falls, reaching zero when they are
        equal.}
  \item Explain why a student who answers ``by Le Ch\^atelier's principle,
        the equilibrium shifts right'' would not receive credit here.
        \answerlines[3]{Because the cell is not at equilibrium --- that is
        precisely why it has a voltage. The CED states that equilibrium
        arguments do not apply to electrochemical systems. The correct
        reasoning compares $Q$ with $K$ and asks whether the change moves
        the cell nearer to or further from equilibrium.}
\end{enumerate}

\begin{exitticket}
A cell's voltage has fallen to zero. State what is true of $Q$, of $K$, and
of $\Delta G$, and whether any reactant remains.
\keyonly{\begin{apcanswer}$Q = K$, so the cell has reached equilibrium, and
$\Delta G = 0$. Reactant almost certainly remains --- equilibrium is not the
same as exhaustion. There is simply no longer any driving force to move
electrons.\end{apcanswer}}
\end{exitticket}

% =====================================================================
\blocklesson{Electrolysis and Faraday's Law}{Zumdahl \S18.7--18.8}

\begin{objectives}
  \item Relate current, time, charge and moles of electrons.
  \item Calculate the mass deposited or dissolved at an electrode.
\end{objectives}

\reading{Zumdahl \S18.7--18.8}{913--925}

\begin{warmup}[3]
  \item Electrolytic cells have $E^\circ$:
        \answerblank[2cm]{negative}
  \item Faraday's constant: \answerblank[3.4cm]{96,485 C/mol e$^-$}
  \item Reduction still occurs at the:
        \answerblank[2.4cm]{cathode}
\end{warmup}

\segment{INSTRUCTION A}{25}{Counting electrons with an ammeter}

\notesectionz{The bridge from electricity to moles}{18.8}
\practice{5}

Electrolysis is stoichiometry in which one of the reagents is
\answerblank[2.4cm]{the electron}. The whole calculation rests on being able
to count electrons, and current is how you do it:

\[ \boxed{\;I = \frac{q}{t}\;} \qquad\text{so}\qquad
   q = It \]
with $q$ in coulombs, $I$ in amperes and $t$ in
\answerblank[1.8cm]{seconds}. Then Faraday's constant converts charge to
moles of electrons:
\[ \text{mol e}^- = \frac{q}{96485} \]

The full chain, worth memorizing as a sequence:
\begin{center}
\fbox{\parbox{0.9\linewidth}{\centering
current \& time $\to$ charge ($q = It$) $\to$ moles of electrons
($\div F$) \\
$\to$ moles of substance ($\div n$ from the half-reaction) $\to$ grams
($\times$ molar mass)}}
\end{center}

\begin{apctrap}
\textbf{The half-reaction supplies the electron ratio, and it is not always
2.} Depositing one mole of metal takes one mole of electrons for
\ce{Ag+}, two for \ce{Cu^2+}, three for \ce{Al^3+}. Reading that number off
the charge of the ion is the step students skip.

Also: \textbf{time must be in seconds}. A problem stated in minutes or hours
is stated that way on purpose.
\end{apctrap}

\begin{workedexample}[Worked example 7: silver plating]
A current of \SI{2.00}{\ampere} runs for \SI{30.0}{\minute} through a
solution of \ce{AgNO3}. What mass of silver plates out?

\textbf{Charge:} $t = 30.0 \times 60 = \SI{1800}{\second}$, so
$q = (2.00)(1800) = \SI{3600}{\coulomb}$.

\textbf{Moles of electrons:}
$3600 / 96485 = \SI{0.0373}{\mole}$.

\textbf{Moles of silver:} the half-reaction is
$\ce{Ag+ + e^- -> Ag}$, one electron per atom, so
$\SI{0.0373}{\mole}$ of \ce{Ag}.

\textbf{Mass:} $(0.0373)(107.87) = \mathbf{4.03}~\text{g}$
\end{workedexample}

\begin{workedexample}[Worked example 8: the same current, a different ion]
Now \SI{1.50}{\ampere} for \SI{1.00}{\hour} through \ce{CuSO4}.

$q = (1.50)(3600) = \SI{5400}{\coulomb}$;
$\text{mol e}^- = 5400/96485 = 0.0560$.

But $\ce{Cu^2+ + 2e^- -> Cu}$ needs \emph{two} electrons per atom:
\[ \text{mol Cu} = \frac{0.0560}{2} = 0.0280
   \qquad m = (0.0280)(63.55) = \mathbf{1.78}~\text{g} \]

Forgetting to divide by 2 gives 3.56 g --- exactly double, and a lost point.
\end{workedexample}

\segment{GUIDED PRACTICE}{15}{Running the chain}

\begin{enumerate}[leftmargin=*,itemsep=0.4em]
  \item How long, in hours, to deposit \SI{5.00}{\gram} of copper at
        \SI{2.00}{\ampere}? \workspace[2.6cm]
        \keyonly{\begin{apcanswer}mol Cu $= 5.00/63.55 = 0.0787$;
        mol e$^- = 0.157$; $q = (0.157)(96485) = \SI{15183}{\coulomb}$;
        $t = 15183/2.00 = \SI{7591}{\second} =
        \mathbf{2.11}~\text{h}$\end{apcanswer}}
  \item What current deposits \SI{10.0}{\gram} of silver in
        \SI{1.00}{\hour}? \workspace[2.4cm]
        \keyonly{\begin{apcanswer}mol Ag $= 10.0/107.87 = 0.0927$;
        mol e$^- = 0.0927$; $q = \SI{8945}{\coulomb}$;
        $I = 8945/3600 = \mathbf{2.48}~\text{A}$\end{apcanswer}}
\end{enumerate}

\segment{INSTRUCTION B}{20}{Where this is used}

\notesectionz{Electrolysis in industry}{18.8}
\practice{2}

\begin{itemize}[leftmargin=*,itemsep=0.3em]
  \item \textbf{Producing aluminium} (the Hall--H\'eroult process).
        \ce{Al^3+} needs three electrons per atom, and aluminium sits at
        $-1.66$~V, so the process is enormously energy-hungry --- which is
        why aluminium smelters are built next to power stations and why
        recycling aluminium saves so much energy.
  \item \textbf{Electroplating} --- depositing a thin layer of one metal
        onto another, controlled precisely by current and time.
  \item \textbf{Electrorefining} --- purifying copper by dissolving the
        impure metal at the anode and depositing pure metal at the cathode.
  \item \textbf{Electrolysis of water}, which has
        $E^\circ = -1.23$~V and so requires an applied potential. This was
        Chapter 17's ``external energy source driving an unfavorable
        process,'' realized in hardware.
\end{itemize}

\segment{APPLICATION}{20}{Scaling up}

\begin{enumerate}[leftmargin=*,itemsep=0.4em]
  \item Chromium is plated from \ce{Cr^3+} at \SI{5.00}{\ampere} for
        \SI{2.00}{\hour}. Find the mass deposited. \workspace[2.8cm]
        \keyonly{\begin{apcanswer}$q = (5.00)(7200) =
        \SI{36000}{\coulomb}$; mol e$^- = 0.373$;
        mol Cr $= 0.373/3 = 0.124$;
        $m = (0.124)(52.00) = \mathbf{6.47}~\text{g}$\end{apcanswer}}
  \item How much charge is needed to produce \SI{1.00}{\kilo\gram} of
        aluminium? \workspace[2.6cm]
        \keyonly{\begin{apcanswer}mol Al $= 1000/26.98 = 37.1$;
        mol e$^- = 3(37.1) = 111$;
        $q = (111)(96485) = \mathbf{1.07\times10^{7}}~\text{C}$ --- about
        ten million coulombs per kilogram\end{apcanswer}}
  \item Two cells are wired in series so the same current passes through
        both: one plates \ce{Ag} from \ce{Ag+}, the other \ce{Cu} from
        \ce{Cu^2+}. After the same time, which has more \emph{moles} of
        metal deposited, and why? \answerlines[3]{The silver cell. The same
        current for the same time delivers the same charge and therefore
        the same moles of electrons to each. Silver needs one electron per
        atom and copper needs two, so twice as many moles of silver are
        deposited.}
\end{enumerate}

\begin{exitticket}
List the four conversion steps between a measured current and a mass of
plated metal, in order.
\keyonly{\begin{apcanswer}Current and time give charge ($q = It$); charge
divided by Faraday's constant gives moles of electrons; moles of electrons
divided by the number of electrons in the half-reaction gives moles of
metal; moles times molar mass gives the mass.\end{apcanswer}}
\end{exitticket}

\end{document}
